\title{LongRoPE: Extending LLM Context Window Beyond 2 Million Tokens}

\begin{abstract}
Expanding the context window in large language models (LLMs) is highly sought after. Nonetheless, present-day context windows remain restricted to approximately 128k tokens, primarily due to the significant expenses associated with fine-tuning, the limited availability of sufficiently long textual data, and the problematic values that arise from previously unseen token positions.

This study presents LongRoPE, marking a significant advancement by expanding the context window of pre-trained LLMs to a remarkable \textbf{2048k} tokens. This is accomplished with as few as 1k fine-tuning steps using training lengths up to 256k, all while preserving the original performance in short context settings. The method centers on three main technical contributions: (i) leveraging and precisely characterizing two distinct non-uniformities present in positional interpolation through an efficient search procedure, which delivers superior fine-tuning initialization and permits an 8$\times$ context increase even without fine-tuning; (ii) deploying a staged extension process whereby a 256k-length LLM is first fine-tuned, followed by a secondary positional interpolation applied to this extended model to enable scaling up to a 2048k context window; (iii) recalibrating LongRoPE at an 8k context length to ensure the restored efficacy for shorter input scenarios. Comprehensive evaluations on LLaMA2 and Mistral models, spanning a variety of tasks, validate the practical benefits of the approach. LongRoPE-enabled models keep their core architectures unchanged, apart from small modifications to positional embeddings, thus supporting compatibility with most existing optimization methods. The implementation will be released at \texttt{https://github.com/microsoft/LongRoPE}.

\section{Introduction}

Large Language Models (LLMs) have achieved significant advancements across a range of applications~\cite{gpt4,llama2}; however, their utility is often hindered by constraints in context window length—for instance, LLaMA2 is restricted to 4096 tokens~\cite{llama2}. When input exceeds this limit, these models experience a degradation in performance, attributable to insufficient exposure during training to extended positional ranges. This limitation becomes particularly problematic in critical use cases such as in-context learning involving numerous demonstration examples~\cite{Huang2023FewerIM} and for LLM-driven agents~\cite{park2023generative,madaan2023selfrefine}.

Recent studies have demonstrated that it is possible to enlarge the context window of a pre-trained LLM to approximately 128k tokens by leveraging fine-tuning on longer sequences~\cite{longlora,pi,yarn,zhang2024soaring,liu2023scaling}. However, there are three significant challenges impeding further increases in context length. Firstly, the introduction of previously unseen position indices causes substantial disruptive values, resulting in distribution shifts and making the fine-tuning process less stable and more difficult to achieve convergence~\cite{pi}. This problem becomes especially pronounced during expansions from 4k to $>$1000k tokens, where over 90\% of position indices are newly incorporated. Secondly, effective fine-tuning necessitates the availability of input sequences matching the target lengths. Present datasets, particularly those with texts surpassing 1000k tokens, are scarce. Additionally, training on such ultra-long inputs demands vast computational resources, requiring excessive time and significant GPU allocation. Thirdly, when the context window is extended to extreme lengths, attention becomes diffused over a multitude of token positions. This broad distribution of attention negatively impacts performance on the shorter contexts seen during pre-training~\cite{pi}.

A commonly employed solution to address the initial challenge involves interpolating RoPE positional embeddings~\cite{rope,pi}, which entails rescaling the position indices to fit within the original pretrained range, as illustrated in Fig.\ref{fig:overview}. Position Interpolation (PI)~\cite{pi} implements a linear interpolation of RoPE's rotary angles based on the extension ratio. NTK~\cite{ntk,dynamicntk} proposes an uneven method that applies distinct interpolation and extrapolation techniques across the various dimensions of RoPE. YaRN~\cite{yarn} further refines this process by segmenting RoPE dimensions into three groups according to frequency, assigning different strategies—extrapolation, NTK, and linear interpolation—to each group respectively. Nonetheless, positional embeddings in the Transformer architecture are characterized by \emph{complex, non-uniform information entropy}. This intricate non-uniformity has not been adequately exploited by current methods, resulting in information loss and thereby constraining the effective context window length.

Section~\ref{sec:analysis} empirically demonstrates two main observations: \textbf{(1)} Robust positional interpolation should account for two distinct non-uniformities—both the disparities in RoPE dimensions and the variability in token positions. In particular, RoPE dimensions at lower values and tokens appearing earlier in the sequence are less sensitive to interpolation; however, the precise approach should be adapted based on the desired length extension. \textbf{(2)} Addressing these sources of non-uniformity within positional interpolation helps preserve the informational structure of the original RoPE, especially with respect to important dimensions and token positions. This approach reduces the detrimental effects of interpolation and thereby yields better model initialization before fine-tuning. Furthermore, it supports extending context length by a factor of 8$\times$ in cases where fine-tuning is not applied.

Driven by these results, we propose \textbf{LongRoPE}, a robust technique designed to increase the LLM context window to over 2 \emph{million} tokens. The development of LongRoPE involves three major contributions. Firstly, LongRoPE leverages the multidimensional variability within positional interpolation to its fullest extent. By pinpointing optimal rescale factors for the rotation angles in each dimension of RoPE, tailored to the location of tokens, the method enhances positional encoding. Given the exponentially growing search space for these rescale factors as the extension ratio rises, LongRoPE incorporates an evolutionary search procedure, utilizing two specialized optimization strategies to streamline the search process. An illustration of the optimized rescaled RoPE can be found in Fig.~\ref{fig:overview}.

Subsequently, {LongRoPE} employs a resource-efficient and incremental extension method to reach a context window of 2048k, circumventing the necessity for fine-tuning on exceptionally long sequences, which are typically rare and difficult to procure. This process initiates by identifying a suitable 256k length for the pre-trained LLM, which is then fine-tuned within this range. Owing to our use of non-uniform positional interpolation—enabling up to an 8$\times$ increase in context window size without further fine-tuning—a subsequent search for optimal RoPE rescale factors is carried out on the already fine-tuned extended LLM. In doing so, we succeed in expanding the context window to 2048k for both LLaMA2 and Mistral~\cite{mistral}.

To prevent losses in performance when dealing with the default, smaller context windows, {LongRoPE} maintains refinement of the RoPE rescale parameters even after the LLM context window has been extended. Analogous to the earlier upscaling process (from 256k to 2048k tokens), we also implement downscaling to 4k and 8k context sizes on a 256k-trained LLM, employing our search method to reduce the degree of positional interpolation. For inference scenarios where the sequence length falls below 8k, RoPE is recalibrated using these identified rescale values.

A broad set of experiments involving multiple LLMs and a range of long-context evaluation tasks validate the efficacy of our approach. Our results indicate that LongRoPE consistently sustains low perplexity across evaluation lengths from 4k to 2048k, attains more than 90\% accuracy in passkey retrieval, and achieves benchmark accuracy similar to existing models within a 4096 token context window. Furthermore, LongRoPE is compatible with any LLM employing RoPE embeddings. We plan to publicly release both our codebase and the LongRoPE-2048k models.

\section{Non-uniformity in Positional Interpolation}
\label{sec:analysis}

\subsection{Preliminary}

Transformer models require explicit positional information, often in the form of position embedding, to represent the order of input tokens. Our work focuses on the RoPE~\cite{rope} position embedding, which is widely used in recent  LLMs. For a token at position index $n$, its corresponding RoPE encoding can be simplified as follows:

\begin{equation}
		\fontsize{7.8}{7.8} \selectfont
	\label{eq:rope}
	\left[ cos(n\theta_0), sin(n\theta_0), cos(n\theta_1), \cdot\cdot\cdot, cos(n\theta_{d/2-1}), sin(n\theta_{d/2-1}) \right]   
\end{equation}

Here, $d$ denotes the embedding size, $n\theta_i$ refers to the rotational angle assigned to the token at position $n$, and $\theta_i=\theta^{-2i/d}$ encodes the frequency of rotation. RoPE employs a default value of $\theta=10000$ as its base.

\textbf{\emph{Context window extension ratio $s$} and positional interpolation}. The parameter $s$ is specified as the proportion between the new context length $L'$ and the initial context length $L$, formally written as $s=\frac{L'}{L}$.

To extend the context window from $L$ to $L'$, current positional interpolation methods suggest downscaling rotation frequencies $\theta_i$ based on extension ratio $s$. Let $\beta=\theta^{2/d}$, and $\lambda$ denote the actual rescale factor related to $s$, we   unify these positional interpolation methods as follows:

\begin{equation}	
	\fontsize{7.5}{7.5} \selectfont
	\label{eq:unified_rope}
	\begin{aligned}
		\left[cos\left(\frac{n}{\lambda(\beta)^0}\right), sin \left(\frac{n}{\lambda(\beta)^0}\right), cos\left(\frac{n}{\lambda(\beta)^1}\right), \cdot\cdot\cdot,  sin \left(\frac{n}{\lambda(\beta)^{d/2-1}}\right) \right]   
	\end{aligned}
\end{equation}

\textbf{Linear positional interpolation (PI)}. PI~\cite{pi} employs a strategy of linearly interpolating positional indices, staying within the boundaries set by pre-training length constraints. When extending the sequence by a factor $s$, each positional rotation angle is uniformly scaled down by $\lambda=s$ across every dimension in RoPE. This uniform reduction compresses the positional representations, which in turn diminishes the model's capability to tell apart tokens situated near one another. Consequently, the effectiveness of PI declines as the extension ratio grows large.

\textbf{NTK-based interpolation and extrapolation}. ~\cite{ntk,dynamicntk} analyze RoPE by considering it as a mechanism for encoding information and utilize the Neural Tangent Kernel (NTK) framework~\cite{jacot2018neural,tancik2020fourier}. To address the issue of position crowding found in PI, they propose redistributing the interpolation constraint across different dimensions of RoPE. In this approach, dimensions associated with higher frequencies are scaled less, while those corresponding to lower frequencies are scaled more, enabling both interpolation and extrapolation in the positional encoding, defined by $\lambda=s^{i}$. The enhanced dynamic NTK~\cite{dynamicntk} further modifies the scaling ratio at each position dynamically, depending on the sequence's current length. In contrast to PI, which requires additional fine-tuning, NTK-informed techniques are capable of enlarging context windows without such fine-tuning, though in practice they typically permit a maximum extension ratio of 4$\times$.

\textbf{YaRN}~\cite{yarn} achieves notable advances in positional interpolation accuracy by segmenting the RoPE dimensions according to their frequency characteristics and then applying distinct interpolation techniques to each group. Dimensions with higher frequencies are subject to extrapolation ($\lambda$=1), whereas those with lower frequencies utilize standard linear interpolation (PI). Intermediate frequency RoPE dimensions are processed via the NTK method. The central innovation of YaRN is its strategy for categorizing RoPE dimensions; however, this process is currently guided by manual, empirical selection, which could hinder performance optimization when adapting to novel LLMs.

\subsection{Study on Non-uniform Positional Interpolation}

Drawing from the approaches of NTK and YaRN, it is evident that their improvements stem from introducing non-linear elements, notably by treating frequencies along RoPE dimensions differently to enhance both interpolation and extrapolation capabilities. Despite this, present implementations of non-linearity predominantly depend on manually crafted heuristics. This leads us to ask: (1) Does existing positional interpolation achieve optimal performance? (2) Are there novel non-linearity mechanisms yet to be discovered?

To address these inquiries, we employ evolution search (refer to Sec.~\ref{sec:method}) to identify improved non-uniform positional interpolation schemes for LLaMA2-7B. Perplexity serves as the search objective, evaluated on five randomly selected samples from the PG19~\cite{pg19} validation dataset. Our empirical study uncovers several important insights as outlined below.

\textbf{Finding 1}: \textit{RoPE dimensions exhibit substantial non-uniformities, which are not effectively handled by current positional interpolation methods.}

We individually optimize the value of $\lambda$ corresponding to each RoPE dimension as specified in Eq.~\ref{eq:unified_rope}. Table~\ref{tbl:exp1} presents a perplexity comparison for LLaMA2-7B across various interpolation methods on the PG19 and Proof-pile~\cite{proof-pile} evaluation sets, all without additional fine-tuning. Our optimized approach delivers notable perplexity reductions, indicating that existing linear (PI) and non-uniform (Dynamic-NTK and YaRN) schemes do not yield optimal results. Interestingly, YaRN performs worse than both PI and NTK on PG19; this is likely due to its inability to fully utilize the intended context window size when applied to LLMs that have not undergone further fine-tuning. As a specific instance, the perplexity for YaRN increases sharply once the context extends past 7k tokens out of an 8k window.

Our investigation revealed that the rescaled factors $\lambda$ in Eq.~\ref{eq:unified_rope} are heterogeneous, contrasting with the constant scaling factor $s$ used in PI, NTK’s analytical approach, and YaRN’s method that employs a group-wise scaling. The adoption of these variable scaling factors yields marked enhancements in LLaMA2’s language modeling capabilities, particularly in terms of perplexity, for context windows of 8k and 16k tokens, all without the need for additional fine-tuning. This improvement is attributed to the adjusted positional embedding, which more faithfully preserves the characteristics of the original RoPE—especially in key embedding dimensions—thereby alleviating the challenge faced by LLMs when differentiating between positions of adjacent tokens.

\textbf{Finding 2}: \textit{RoPE for the initial tokens in the input sequence should be extrapolated with less interpolation.} 

We propose that for the first $\hat{n}$ tokens within input sequences, the RoPE mechanism should perform minimal interpolation. This is grounded in the fact that these tokens typically attract higher attention scores, thereby serving as pivotal elements within attention modules, as previously demonstrated by Streaming LLM~\cite{streamingllm} and LM-Infinite~\cite{han2023lminfinite}. To test this conjecture, we scale up the context window to 8k and 16k tokens utilizing PI and NTK approaches, while omitting interpolation for the initial  $\hat n$ tokens (with values 0, 2, ..., 256). When $\hat n$ equals zero, this returns to the baseline PI and NTK methods. As summarized in Table~\ref{tbl:tokenposition}, two important insights emerge: \textbf{(1)} exempting the earliest tokens from position interpolation yields tangible gains in model performance; \textbf{(2)} the optimal choice of $\hat n$ is influenced by the specific context window length employed.

\textbf{Finding 3}: \textit{Non-uniform positional interpolation effectively extends LLM context window in both fine-tuning and non-fine-tuning settings.} 

Although our findings demonstrate that the proposed non-uniform position interpolation yields notable gains in extension performance at the 8k and 16k context sizes without additional fine-tuning, extending beyond these ranges necessitates further model adjustment. To this end, we conduct fine-tuning of LLaMA2-7B with our optimized RoPE scheme tailored for a 64k context window (see Appendix for experimental configurations). The results presented in Table~\ref{tbl:finetune} indicate that our approach achieves substantial performance improvements over both PI and YaRN, across both the pre- and post-fine-tuning stages of LLaMA2-7B. This advantage stems from the refined application of non-uniform positional interpolation, which reduces information degradation and supplies a stronger foundation for effective fine-tuning.

\textbf{Summary}. This research identifies two sources of non-uniformity: variability across RoPE dimensions and token position indices. Effectively leveraging these non-uniform features in positional interpolation substantially enhances the ability of LLMs to handle longer contexts.

\section{LongRoPE}

\label{sec:method}
Building on these insights, we propose LongRoPE. This approach initially employs an efficient search algorithm specifically designed to leverage both non-uniformities, and subsequently utilizes this strategy to push the LLM context window to lengths exceeding 2 million tokens.

\subsection{Problem Formulation}

These two forms of non-uniformity may expand the solution space substantially and complicate the optimization process. To tackle this challenge, we recast the multidimensional non-uniform position interpolation optimization issue as a search problem.

For a LLM targeting a  context window size of $L'$ and lengthy input documents $\mathbf{X}$, where each $\mathbf{x}\in\mathbf{X}$ surpasses $L'$ in token length, we denote the original rotary angle of the $i^{th}$ dimension in RoPE embedding at token position $n$ as  $\frac{n}{\beta^i}$. The optimization problem is then formulated as follows:

\begin{equation}
\begin{gathered}
		\label{eq:problem}

\underset {\mathbf{x}\in\mathbf{X}; \, |\mathbf{x}|\ge L'}{ \operatorname {arg\,min} } \,  \mathcal{L}\, \left( \text{LLM} (\text{RoPE}, \mathbf{X})\right),\ \text{with} \
\\
\scalebox{0.93}{
$\underset {\begin{subarray}{l} i=0,\ldots,\frac{d}{2}-1;\\ n\in[ 0,|\mathbf{x}|);\end{subarray}}{\text{RoPE}_(n)}= {\left[\ldots,\,\cos\left(\mathbb{I}(\hat{\lambda}_{i}, \hat{n})\frac{n}{\beta^i}\right),\, \sin\left(\mathbb{I}(\hat{\lambda}_{i}, \hat{n})\frac{n}{\beta^i}\right),\,\ldots\right] }$}\\
\text{where} \quad \mathbb{I}(\hat{\lambda}_{i},\hat{n}) = \begin{cases}
1 & \text{if}\ n < \hat{n} \\
{\frac{1}{\lambda}_{i}} & \text{if}\ n \geq \hat{n}
\end{cases}
	\end{gathered}
\end{equation}

In this context, we define a set of scaling coefficients, $\mathbb{I}(\hat{\lambda}_{i},\hat{n})$, to account for the presence of two distinct types of non-uniformity. Here, $\hat{\lambda}_i$ reflects dimension-specific non-uniformity within RoPE, and $\hat{n}$ specifies positional irregularities among tokens. The purpose of $\mathbb{I}(\hat{\lambda}_{i},\hat{n})$ is to adjust the rotation angle for the $i^{th}$ RoPE dimension, where $\hat\lambda_{i}$ functions as the scaling coefficient and $\hat{n}$ marks a positional threshold. For the first $\hat{n}$-1 token positions, the original RoPE rotary angle $\frac{n}{\beta^i}$ remains unchanged, since the scaling factor $\hat\lambda_{i}$ is not applied. When tokens occur at positions $n\ge\hat{n}$, the scaling factor $\hat\lambda_{i}$ is introduced.

Assuming a desired context window of length $L'$, our task is to determine the set of optimal rescale factors ($\mathbb{I}(\hat{\lambda}_{0},\hat{n})$, $\mathbb{I}(\hat{\lambda}_{1},\hat{n})$, ..., $\mathbb{I}(\hat{\lambda}_{i},\hat{n})$, ...) across each RoPE dimension from the first through the $d^{th}$. Implementing these rescaled factors in the target $\text{LLM}$, we aim to minimize the next token prediction loss $\mathcal{L}$ (equivalent to perplexity) when processing input sequences $\mathbf{X}$ of length $L'$.

\subsection{Searching the Non-uniform Position Interpolation}

To address the challenge outlined in Eq.~\ref{eq:problem}, we propose an efficient and straightforward approach that identifies optimal RoPE rescaling factors, thereby taking full advantage of the non-uniformities present across multiple dimensions within the position embedding.

\textbf{Search space}. We construct an extensive search space that accommodates both forms of non-uniformity. The structure of this search space is detailed in Table~\ref{tbl:searchspace}. In particular, the approach enables the selection of a unique rescale factor for each individual dimension in the RoPE embedding. To streamline the search process, we optimize over $\lambda_i$ and $\hat{n}$ instead of directly exploring $\mathbb{I}(\hat{\lambda}_{i},\hat{n})$, with the relationship defined as $\hat{\lambda}_{i}=1/\lambda_i$. Table~\ref{tbl:searchspace} indicates that $\lambda_i$ can be searched within the range starting from 1.0 (representing pure extrapolation) up to $s\times1.25$ (corresponding to a broader interpolation than PI), incremented by 0.01 per step. Here, $s$ denotes the ratio used to extend the desired context window.

The parameter $\hat{n}$ determines how many of the first token positions preserve their original RoPE embeddings without undergoing position interpolation. In practice, $\hat{n}$ is selected from the set \{0, 1, 2, 4, 8, 12, 16, 20, 24, 28, 32, 64, 128, 256\}. If $\hat{n}=0$, then every token position applies the rescale factors identified during the search process.

\textbf{Evolution-based search}. The search space outlined in Table~\ref{tbl:searchspace} encompasses a wide array of positional interpolation options, making exhaustive exploration highly inefficient. For instance, expanding by $s=4\times$ results in $400^{128/2}\times14=4\times10^{167}$ possible configurations. As the extension ratio increases, the size of the search space grows at an exponential rate. To efficiently navigate this vast landscape, we employ evolution search~\cite{guo2020single}, supplemented with two key optimization methods designed to significantly enhance the effectiveness of our search process. The overall approach is detailed in Algorithm~\ref{alg:search}.

\textit{Optimized initial population generation}. Rather than populating the initial group of $P$ rescale factors purely at random, we incorporate the specific RoPE rescale factors associated with PI, NTK, and YaRN directly as distinct individuals within the initial population. The rest of the $P$-3 members are produced by applying random mutations to these three rescale factors, each mutation occurring with probability $p$.

\textit{Monotonically non-decreasing constraint}. Following the initialization of the population, LLM perplexity is evaluated for every candidate. In detail, the designated RoPE rescale factors are applied within the target LLM, and perplexity for the input $\mathbf{X}$ is subsequently measured. The individuals with the lowest perplexity scores are selected as the top-k parents to guide the evolutionary process. Nonetheless, the immense breadth of the search space means that simple mutation and crossover operations may direct the search toward suboptimal candidates, resulting in excessive perplexity evaluations. This inefficiency is especially pronounced when $L'$ is large, since each perplexity calculation requires substantial inference time.

To tackle this issue, we enforce a monotonicity condition such that the rescaled RoPE factors sampled must not decrease, formally $\lambda_i \leq \lambda_{i+1}$. Only those RoPE configurations adhering to this rule are used in the large language model for perplexity assessments, resulting in a substantial reduction in computational overhead. Concretely, the constraint stipulates that $\lambda_i$ should progress monotonically with respect to the RoPE dimension index ($i = 0,\ldots, 63$). This requirement stems from the Neural Tangent Kernel (NTK) framework~\cite{jacot2018neural,tancik2020fourier,ntk}, which posits that the dimensions corresponding to higher frequencies (i.e., lower $i$) warrant less interpolation—implying a smaller value for $\lambda_i$—whereas those associated with lower frequencies (i.e., higher $i$) permit greater interpolation, thus justifying larger $\lambda_i$.

\begin{algorithm}[t]
	\small
	\caption{The search algorithm}
	\textbf{Input:} target LLM, input samples $\mathbf{X}$, population size $P$, mutation size $N_1$, crossover size $N_2$, maximum iterations $\mathcal{T}$, mutation probability $p$\\

	\begin{algorithmic}[1]
		\label{alg:search}
		\STATE Initialize Top-k as empty;
		\STATE Set $\text{P}_0$ to the result of \textit{Initialize\_population\_with\_optimization} ($P$, $\mathbf{X}$, $p$);
		\FOR{$i$ = $1$ to $\mathcal{T}$}
		\STATE Evaluate perplexity for $\text{P}_{i-1}$ over $\mathbf{X}$ using the LLM;
		\STATE Update Top-k by invoking \textit{Update\_Topk} with the current population;
		\STATE Generate $\text{P}_{mutation}$ using \textit{Mutation\_with\_mono\_constraint} on Top-k, with size $N_1$ and probability $p$;
		\STATE Obtain $\text{P}_{crossover}$ by applying \textit{Crossover\_with\_mono\_constraint} on Top-k and size $N_2$;
		\STATE Merge $\text{P}_{mutation}$, $\text{P}_{crossover}$, and Top-k to form $\text{P}_i$;
		\ENDFOR
		\STATE Output the member of Top-k possessing the minimal perplexity;
	\end{algorithmic}
\end{algorithm}

\textbf{8$\times$ extension without fine-tuning}. Through evolutionary search, our approach efficiently discovers non-uniform RoPE rescale parameters that safeguard crucial dimensions and positions, thereby reducing information degradation from interpolation. As illustrated in Fig.\ref{fig:non-ft}, this technique allows the LLaMA2 context window to be expanded from 4k to 32k without requiring any fine-tuning. By comparison, alternative approaches—such as PI, as well as non-uniform NTK and YaRN—exhibit a significant increase in perplexity after just a 2$\times$ window extension.

\subsection{Extending LLM Context Window to 2048K}
\label{sec:extendto2048k}

\textbf{Progressive extension to 2048k}. In this section, we present our approach for scaling the context window of pre-trained LLMs from the typical 4k up to more than 2048k tokens. Our non-uniform positional interpolation allows for an 8$\times$ increase without the need for fine-tuning. However, when aiming for substantially greater extensions (such as 512$\times$), fine-tuning becomes indispensable. One strategy involves identifying appropriate RoPE rescaling factors for the target 2048k window followed by a fine-tuning process. This approach, though, is hampered by the immense computational cost required for training at such a large scale. Additionally, we have observed that it is difficult to fine-tune LLMs effectively for large extension ratios (refer to the Appendix for further discussion).

LongRoPE demonstrates efficacy for both the base model and its fine-tuned extended variants. Consequently, we propose a progressive and efficient strategy that reaches the intended 2048k context length using only 1,000 fine-tuning steps, all conducted at a training sequence length capped at 256k.

$\Diamond$ \textit{LongRoPE-based search for scaling up pre-trained LLMs to 256k.} Using LLaMA2 as a reference model, we perform a search process aimed at increasing the context window to 128k and 256k tokens. At this phase, the respective extension ratios achieved are 32$\times$ and 64$\times$.

$\Diamond$ \textit{Fine-tuning to 256k}. Subsequently, we adapt the pre-trained LLM to handle a context window of 256k tokens. In practice, our approach begins by fine-tuning LLaMA2 over 400 iterations, utilizing RoPE rescaling parameters set for 128k. After this stage, we update the RoPE rescaling factors to support 256k on the resulting checkpoint and proceed with 600 further fine-tuning steps. This sequential strategy demonstrates greater efficiency compared to a single-stage fine-tuning directly towards 256k.

$\Diamond$ \textit{LongRoPE search enables 2048k extension for fine-tuned extended LLM}. In the last step, we conduct an additional search operation using the previously fine-tuned LLM with a context length of 256k. As a result, we are able to achieve a dramatically expanded context window of 2048k, eliminating the need for any further fine-tuning. This yields a total extension factor of $512\times$.

\textbf{Degradation within original context window}.  When scaling up to a notably large 2048k context window, a reduction in performance is observed for inputs falling inside the initial context window range. This phenomenon is a recognized limitation of positional interpolation~\cite{pi}, as it compresses the positional embeddings for the earlier segments of the sequence into a relatively tight space in higher-dimensional representations. This crowding of embeddings impairs the model’s effectiveness. Specifically, with an extension factor of 512$\times$, the embeddings for positions in the standard 4k context window experience significant overlap.

To address this issue, we conduct an additional evolution search on the expanded LLM, specifically tuning RoPE rescale parameters for shorter context lengths such as 4k and 8k. For these shorter sequences, we limit the maximum $\lambda$ considered, since less positional interpolation is necessary. At inference time, the LLM adapts the relevant RoPE rescale factors in response to the context length.

\section{LongRoPE}

\label{sec:method}
Building on these observations, we propose LongRoPE. This approach begins by leveraging a novel and efficient search algorithm designed to capitalize on the identified dual non-uniformities. Subsequently, this algorithm enables the extension of the LLM context window to support more than 2 million tokens.

\subsection{Problem Formulation}

These two forms of non-uniformity expand the range of possible solutions and complicate the optimization process. To manage this challenge, we recast the problem of optimizing multidimensional non-uniform position interpolation as a search problem.

For a LLM targeting a  context window size of $L'$ and lengthy input documents $\mathbf{X}$, where each $\mathbf{x}\in\mathbf{X}$ surpasses $L'$ in token length, we denote the original rotary angle of the $i^{th}$ dimension in RoPE embedding at token position $n$ as  $\frac{n}{\beta^i}$. The optimization problem is then formulated as follows:

\begin{equation}
\begin{gathered}
		\label{eq:problem}

\underset {\mathbf{x}\in\mathbf{X}; \, |\mathbf{x}|\ge L'}{ \operatorname {arg\,min} } \,  \mathcal{L}\, \left( \text{LLM} (\text{RoPE}, \mathbf{X})\right),	\text{where}
\\
\scalebox{0.93}{
$\underset {\begin{subarray}{l} i=0,\cdot\cdot,\frac{d}{2}-1;\\ n\in[ 0,|\mathbf{x}|);\end{subarray}}{\text{RoPE}_(n)}= {\left[\cdots,\cos\left(\mathbb{I}(\hat{\lambda}_{i}, \hat{n})\frac{n}{\beta^i}\right), \sin\left(\mathbb{I}(\hat{\lambda}_{i}, \hat{n})\frac{n}{\beta^i}\right),\cdots\right] }$}\\
	\text{where} \quad \mathbb{I}(\hat{\lambda}_{i},\hat{n})=\begin{cases}
	1&\text{if $n<\hat{n}$}\\
{\frac{1}{\lambda}_{i}}&\text{if $n\geq\hat{n}$}
\end{cases}
	\end{gathered}
\end{equation}

We propose a collection of rescaling coefficients, $\mathbb{I}(\hat{\lambda}_{i},\hat{n})$, which are designed to account for two distinct types of non-uniformity. The parameters $\hat{\lambda_i}$ and $\hat{n}$ correspond to non-uniformities in RoPE dimension scaling and token positional differences, respectively. In detail, $\mathbb{I}(\hat{\lambda}_{i},\hat{n})$ is employed to adjust the rotation angle associated with the $i^{th}$ dimension of RoPE, using $\hat\lambda_{i}$ as the rescaling coefficient and $\hat{n}$ as the threshold for token positions. For the first $\hat{n}-1$ token positions, the standard RoPE rotary angle $\frac{n}{\beta^i}$ is maintained, unaffected by $\hat\lambda_{i}$. When the token position reaches $n \ge \hat{n}$, the rescaling coefficient is activated and modifies the angle accordingly.

Suppose a desired context window length of $L'$ is specified; the aim is to determine the set of optimal rescaling coefficients ($\mathbb{I}(\hat{\lambda}_{0},\hat{n})$, $\mathbb{I}(\hat{\lambda}_{1},\hat{n})$, ...$\mathbb{I}(\hat{\lambda}_{i},\hat{n})$...) corresponding to the $1^{st}$ through $d^{th}$ dimensions of RoPE. Implementing these rescale factors in the target $\text{LLM}$'s $\text{RoPE}$ is intended to minimize the next-token prediction loss, denoted $\mathcal{L}$ (which equates to perplexity), across input sequences $\mathbf{X}$ comprising $L'$ tokens.

\subsection{Searching the Non-uniform Position Interpolation}

To address the issue presented in Eq.~\ref{eq:problem}, we propose an efficient and straightforward approach that identifies the most suitable RoPE rescale parameters, thereby maximizing the utilization of multidimensional irregularities inherent in position embeddings.

\textbf{Search space}. A comprehensive search space is constructed to incorporate both forms of non-uniformity. The search space is detailed in Table~\ref{tbl:searchspace}. In particular, the approach enables the selection of an individualized rescale factor for each RoPE embedding dimension. For greater simplicity, the search is performed over $\lambda_i$ and $\hat{n}$, as opposed to directly optimizing $\mathbb{I}(\hat{\lambda}_{i},\hat{n})$, where  $\hat{\lambda}_{i}=1/\lambda_i$. As indicated in Table~\ref{tbl:searchspace}, $\lambda_i$ is permitted to vary within the range starting from 1.0 (representing straightforward extrapolation) up to $s\times1.25$ (corresponding to increased interpolation compared to PI) in increments of 0.01. Here, $s$ denotes the ratio by which the context window is extended.

The parameter $\hat{n}$ specifies how many starting token positions maintain the unaltered RoPE embedding, foregoing positional interpolation. In practice, we select $\hat{n}$ from the set \{0, 1, 2, 4, 8, 12, 16, 20, 24, 28, 32, 64, 128, 256\}. If $\hat{n}=0$, then every token position utilizes the optimized rescale factors.

\textbf{Evolution-based search}. The search space outlined in Table~\ref{tbl:searchspace} encompasses a wide array of positional interpolation strategies, resulting in substantial complexity when it comes to systematic exploration. For instance, a scaling factor of $s=4\times$ generates $400^{128/2}\times14 = 4\times10^{167}$ possible configurations. As the extension ratio increases, the search space grows at an exponential rate. To navigate this, we employ evolution search~\cite{guo2020single}, and implement two optimization methods to significantly enhance the efficiency of the search process. The full search methodology is depicted in Algorithm~\ref{alg:search}.

\textit{Refined strategy for initial population}. Rather than relying solely on random selection to create an initial population of $P$ rescale factors, we explicitly insert the three RoPE rescale factors that correspond to PI, NTK, and YaRN into the initial group. The additional $P$-3 population members are generated by introducing random mutations to these three rescale factors with mutation occurring at probability $p$.

\textit{Monotonically non-decreasing constraint}. Following the creation of the initial population, we evaluate each candidate by calculating LLM perplexity. For this, we assign the specific RoPE rescale factors to the target LLM and determine the perplexity over input $\mathbf{X}$. The top-k candidates with the lowest perplexity scores are selected as parents for subsequent evolutionary steps. Nevertheless, the expansive nature of the search space means that standard mutation and crossover may frequently result in suboptimal candidates, incurring redundant perplexity computations. This issue becomes especially pronounced for large $L'$, since each perplexity evaluation is computationally intensive.

To mitigate this issue, we enforce a non-decreasing monotonicity condition on the rescaled RoPE factors sampled: $\lambda_i \le \lambda_{i+1}$. Only those RoPE settings meeting this requirement are used in perplexity tests on LLMs, which greatly curtails the computational burden associated with searching. More precisely, we stipulate that $\lambda_i$ must rise monotonically as the RoPE dimension increases (where $i = 0, \dots, 63$). The rationale behind this dimensional monotonicity draws upon NTK theory~\cite{jacot2018neural,tancik2020fourier,ntk}, which indicates that lower-indexed dimensions associated with higher frequencies should use less interpolation (smaller values of $\lambda_i$), whereas higher-indexed dimensions tied to lower frequencies allow for greater interpolation (larger $\lambda_i$).

\begin{algorithm}[t]
	\small
	\caption{The search algorithm}
	\textbf{Input:} target LLM, input samples $\mathbf{X}$, population size $P$, mutation size $N_1$, crossover size $N_2$, maximum iterations $\mathcal{T}$, mutation probability $p$\\

	\begin{algorithmic}[1]
		\label{alg:search}
		\STATE Top-k$=\phi$;	
		\STATE $\text{P}_0$=\textit{Initialize\_population\_with\_optimization}($P$, $\mathbf{X}$, $p$); 
		\FOR{$i$=1\ to\ $\mathcal{T}$}
		\STATE \textit{Compute\_perplexity}(LLM, $\text{P}_{i-1}$, $\mathbf{X}$);
		\STATE Top-k = \textit{Update\_Topk}(Top-k, $\text{P}_{i-1}$);
		\STATE $\text{P}_{mutation}$ = \textit{Mutation\_with\_mono\_constraint}(Top-k, $N_1$, $p$);
		\STATE $\text{P}_{crossover}$ = \textit{Crossover\_with\_mono\_constraint}(Top-k, $N_2$);
		\STATE $\text{P}_i$ = $\text{P}_{mutation}$ $\cup$ $\text{P}_{crossover}$ $\cup$ Top-k;
		\ENDFOR
		\STATE Output the candidate in Top-k with minimal perplexity;
	\end{algorithmic}
\end{algorithm}

\textbf{8$\times$ extension without fine-tuning}. Through evolutionary search, our approach efficiently discovers optimal non-uniform RoPE rescaling parameters, maintaining crucial dimensions and positions to reduce the information degradation commonly caused by interpolation. As illustrated in Fig.\ref{fig:non-ft}, this technique enables LLaMA2's context window to expand from 4k to 32k tokens without requiring any fine-tuning. Conversely, prior methods—including PI, as well as non-uniform NTK and YaRN—experience a significant increase in perplexity beyond a 2$\times$ context length extension.

\subsection{Extending LLM Context Window to 2048K}
\label{sec:extendto2048k}

\textbf{Progressive extension to 2048k}. In this section, we present our strategy for expanding the context window of pre-trained LLMs from the conventional 4k limit up to 2048k tokens. Our results indicate that non-uniform positional interpolation facilitates an extension of up to 8$\times$ the original window size without requiring fine-tuning. However, to achieve even greater expansion ratios—such as 512$\times$—fine-tuning becomes indispensable. One potential approach involves identifying suitable RoPE rescaling parameters for the desired 2048k window before proceeding with fine-tuning. Despite this, the process is hampered by the significant computational demands associated with training at such scales. Furthermore, our observations suggest that it remains difficult to effectively fine-tune LLMs under these large extension conditions (refer to Appendix).

LongRoPE demonstrates efficacy when applied to both the base and fine-tuned extended LLM. Thus, we propose a resource-efficient, stepwise approach that reaches the intended 2048k context window using only 1k fine-tuning iterations, all conducted with a maximum training length of 256k.

$\Diamond$ \textit{LongRoPE search for scaling pre-trained LLMs to 256k}. Using LLaMA2 as a case study, we explore context window extension through a search process for 128k and 256k target sequence lengths. At this stage, this corresponds to expansion ratios of 32$\times$ and 64$\times$, in order.

$\Diamond$ \textit{Fine-tuning to 256k}. Next, the pre-trained LLM undergoes fine-tuning to extend its context window to 256k tokens. Initially, we conduct 400 fine-tuning steps on LLaMA2 employing RoPE rescaled factors for 128k. Following this stage, we update the RoPE rescaled factors to support 256k and resume training from the resulting checkpoint for an additional 600 steps. This staged approach is found to be more effective than attempting fine-tuning to 256k in a single phase.

$\Diamond$ \textit{Applying LongRoPE search to expand fine-tuned LLMs to 2048k.} In the final step, we conduct an additional search procedure using the previously fine-tuned LLM at a 256k context length. This enables us to achieve an exceptionally large context window of 2048k tokens without the need for additional fine-tuning. Consequently, the context extension reaches a ratio of $512\times$.

\textbf{Original context window degradation}. Upon scaling up to a substantial 2048k context window, we observe diminished performance in the model's native context span. This decline is well-documented in the literature on positional interpolation~\cite{pi}, which highlights that compressing positional embeddings into a restricted subspace within the base context window can hinder language modeling effectiveness. Specifically, when the extension ratio is as high as 512$\times$, positional representations in the initial 4k context window become highly congested.

To address this issue, we conduct an additional evolution search on the extended LLM, specifically tuning the RoPE rescale factors for shorter context lengths such as 4k and 8k. For these reduced lengths, we constrain the upper limit of the search space for $\lambda$, given the diminished necessity for positional interpolation. At inference time, the LLM adapts the RoPE rescale factors dynamically based on the context length used.

 \section{Experiments}

\label{sec:exp}
\subsection{Setup}
\textbf{Evaluation Tasks and models}. Our investigation employs LongRoPE with both LLaMA2-7B and Mistral-7B architectures, assessing their capabilities along three dimensions: (1) measuring perplexity when processing lengthy documents using extended-context LLMs; (2) performance on the Passkey retrieval task, designed to test the model's aptitude at identifying a straightforward passkey amidst predominantly extraneous content; and (3) outcomes on conventional LLM benchmark tasks restricted to a context window of 4096 tokens.

\textbf{Fine-tuning}. For LLaMA2, the training procedure employs a linear learning rate schedule starting at 2e-5 and a global batch size set to 32. Initially, the model is fine-tuned for 400 iterations on the Redpajama~\cite{redpajama} dataset, partitioned into segments of 128k tokens, each framed with BOS and EOS markers. Subsequently, using the resulting checkpoint, an additional 600 steps are performed to extend the context window to 256k. Fine-tuning with a 128k context size is conducted across 8 A100 GPUs utilizing the distributed training infrastructure~\cite{cube}, while expanding to 256k context involves 16 A100 GPUs. For the Mistral model, a fixed learning rate of 1e-6 is applied, coupled with a global batch size of 64. For both 128k and 256k models, the YaRN~\cite{yarn} configuration is adopted, running 400 steps on the Together Computer's Long-Data Collections~\cite{mistraldata} with sequences capped at 16k tokens. Training is accomplished on 4 A100 GPUs.

\textbf{Search}. For window sizes up to 256k, the following parameters are used: $P$=64, $N_1$=$N_2$=16, $p$=0.3, $\mathcal{T}$=40, with the top 32 candidates selected for mutation or crossover during each cycle. Perplexity evaluation relies on 5 randomly chosen samples from the PG19 validation set, ensuring that each meets at least the target context length. For window sizes above 512k, we reduce the population size as well as the numbers used in mutation and crossover procedures by half. Perplexity in this case is assessed using 3 random samples taken from the Pile-Books3~\cite{pile} validation set.

\textbf{Baselines}. In order to achieve a context length of 2048k, we conducted fine-tuning on models previously trained with 128k and 256k context windows. This process resulted in LongRoPE-2048k (ft=128k) for LLaMA2 and LongRoPE-2048k (ft=256k) for Mistral. These four models are evaluated alongside leading approaches for expanding context windows, focusing particularly on open-source LLMs that have been fine-tuned subsequent to applying positional interpolation through PI, NTK, and YaRN methodologies. The comparison set includes Together-32k~\cite{together}, Code LLaMA~\cite{codellama}, LongLoRA-full-FT-100k~\cite{longlora}, as well as YaRN-LLaMA and YaRN-Mistral~\cite{yarn}.

\subsection{Main Results}

\label{sec:mainresult}
\textbf{Language modeling over long sequences up to 256k tokens}. Our initial experiments benchmark our model against leading extended LLMs, restricting the evaluation context to 256k tokens. To illustrate the robustness of our approach, we employ two distinct datasets: the test portions of Proof-pile~\cite{pg19} and PG19~\cite{pile}. Perplexity is measured at different context lengths, employing a sliding window approach with a window size of 256. For the PG19 benchmark, all 100 documents from the test split are utilized. In the case of Proof-pile, following the procedure of YaRN~\cite{yarn}, we draw 10 random samples, ensuring each sample contains at least 128k tokens.

Table~\ref{tbl:proof-pile} and Table~\ref{tbl:pg19} present a perplexity comparison between LLaMA2 and Mistral models extended through various interpolation strategies on the Proof-pile and PG19 benchmarks, respectively. Two main findings emerge: \textbf{(1)} Our extended models consistently demonstrate reduced perplexity across evaluation lengths from 4k up to 256k, indicating improved ability to utilize extended contexts. \textbf{(2)} Remarkably, despite operating with context windows up to 16$\times$ longer—an environment in which maintaining performance at shorter lengths is generally difficult—our LongRoPE-2048k models achieve superior results relative to leading baselines within the 256k context range.

\textbf{Language modeling for sequences longer than 2000k}. To assess performance on very long textual data, we employ the Books3~\cite{pile} dataset. For practical evaluation, 20 books are randomly chosen, each with a length greater than 2048k, and a sliding window of 256k is applied.

Table~\ref{tbl:books3} demonstrates that LongRoPE effectively expands the context window of both LLaMA2-7B and Mistral-7B models up to 2048k tokens, while maintaining perplexity levels that are on par with or even better than those of the baseline models across shorter windows of 8k to 128k tokens. Additionally, we note substantial performance disparities between the 2048k variants of LLaMA2 and Mistral. In the short-length regime, Mistral achieves superior results compared to baselines, but its perplexity surpasses 7 once the context window goes beyond 256k tokens. LLaMA2's results behave as expected, with a gradual reduction in perplexity as context length increases, albeit with slight upticks at 1024k and 2048k. Notably, for LLaMA2, fine-tuning LongRoPE-2048k at a 256k context length yields improved outcomes relative to fine-tuning at 128k, attributed to the smaller secondary extension ratio (specifically, 8$\times$ as opposed to 16$\times$). Conversely, Mistral attains its best performance when fine-tuned at a window size of 128k. This primarily arises from the use of a 16k training length, following YaRN's configuration, during Mistral's fine-tuning at 128k and 256k windows, which limits the extent to which Mistral’s context window can be expanded post fine-tuning.

\textbf{Passkey retrieval}. In this section, we examine how the context window size impacts generation tasks by utilizing a synthetic passkey retrieval benchmark introduced by~\cite{passkey}. For this evaluation, the model is required to identify a concealed random five-digit passkey within an extended text. Details about the prompt format are provided in the appendix. We conduct the passkey retrieval task over 10 trials, with the passkey’s position randomly assigned, distributed uniformly throughout the evaluation context.

Fig.~\ref{fig:passkey} illustrates retrieval accuracy across different approaches. The performance of current baseline models declines sharply to zero once the token count exceeds 128k. By contrast, LongRoPE-LLaMA2-2048k (ft=256k) achieves consistently high retrieval accuracy ($\ge$90\%) throughout the range from 4k to 2048k tokens, even under the demanding scenario of extracting a passkey from over a million tokens. LongRoPE-Mistral-2048k (ft=128k) sustains perfect accuracy up to 1800k tokens and then decreases to 60\% at 2048k. This trend is consistent with Table~\ref{tbl:books3}, which shows a modest increase in perplexity at 2048k tokens.

\textbf{Assessment on established benchmarks within the original context window}. The performance of LongRoPE-2048k models is assessed on the native context window using the Hugging Face Open LLM Leaderboard~\cite{huggingfaceboard}, examining both zero-shot and few-shot scenarios. Our evaluation protocols include 25-shot for ARC-Challenge~\cite{allenai:arc}, 10-shot for HellaSwag~\cite{zellers2019hellaswag}, 5-shot for MMLU~\cite{mmlu}, and a zero-shot setting for TruthfulQA~\cite{lin2021truthfulqa}.

According to Table~\ref{tbl:commonsense}, our models perform on par with the original benchmark, which was intended for use with shorter context windows, and surpass the original Mistral on TruthfulQA by a margin of +0.5\%. Although LongRoPE-LLaMA2-2048k, trained with a context length of 256k, experiences a modest drop in performance, it still maintains satisfactory results across the majority of tasks.

\subsection{Ablation Results}

\textbf{Assessing the second positional interpolation's impact}. Our progressive extension methodology involves applying a second round of non-uniform positional interpolation, leveraging our search algorithm on the previously fine-tuned and extended LLMs. To assess its utility, we experimented with a fine-tuned LLaMA2-256k model, scaling its context window to 512k, 1024k, and 2048k contexts using both PI and YaRN. Results in Table~\ref{tbl:secondsearch} demonstrate that our non-uniform positional interpolation method maintains stable perplexity metrics. On the other hand, models extended using PI and YaRN experience a rapid increase in perplexity as the context window grows.

\textbf{Effectiveness of recovery at shorter context lengths.} In order to address diminished performance at reduced context lengths, we revise the RoPE scaling parameters for LongRoPE-2048k using our search procedure. Notably, we limit the upper bound of scale factors during the search process to minimize interpolation for sequences of length 4k and 8k. Table~\ref{tbl:shortperformance} presents a comparison of perplexity for LongRoPE-LLaMA2-2048k evaluated on Proof-pile at 4k and 8k context lengths, alongside mean benchmark accuracy for LLMs. The findings highlight substantial gains in performance for these shorter contexts.

\textbf{Analysis of the two types of non-uniformity}. To assess the individual impact of each non-uniformity, we conducted ablation studies. Specifically, we performed two sets of experiments: (i) expanding LLaMA2-7B's context to 16k and 32k tokens using several approaches—PI, optimizing only the RoPE dimension, and optimizing both forms of non-uniformity; (ii) increasing the sequence length for our 256k-length fine-tuned LLaMA2 model to 2048k with the same methodology. Perplexity was measured without further fine-tuning. According to Table~\ref{tbl:searchdimension}, adjusting the RoPE dimension for non-uniformity considerably lowers perplexity relative to the linear interpolation employed in PI. Incorporating non-uniformity in token position yields noticeable performance gains at 16k and 32k, whereas it has minimal effect at 2048k, likely due to the substantial sequence length. Simply retaining initial tokens without any interpolation offers little benefit in this setting; a more detailed investigation is deferred to future research.

\section{Related Works}

Beyond techniques grounded in position interpolation, this section examines alternative methods found in the literature.

\textbf{Retrieval-based} techniques utilize external memory systems to store extended historical context and employ retrieval mechanisms to fetch pertinent documents during inference~\cite{longllama,wang2023augmenting,borgeaud2022improving}. Such strategies often require significant alterations to the underlying LLM architectures. Conversely, our approach is considerably lighter in implementation, only involving minimal adjustments to positional embeddings. Additionally, our method is versatile and applicable to a wider range of long-context tasks beyond simple retrieval, including extended document summarization and few-shot learning scenarios.

\textbf{Extensions of context windows using attention}. In addition to positional embedding interpolation, certain studies have managed to expand the input context by retaining the initial LLM context window size, specifically by altering attention mechanisms~\cite{han2023lminfinite, streamingllm,parallelwindow}. The central strategy involves alleviating the increase in attention complexity triggered by incorporating new positions, accomplished through innovative attention masking approaches. Approaches based on attention masking and those employing positional interpolation are synergistic.

\textbf{Fine-tuning based approaches} investigate strategies to adapt pre-trained LLMs with altered position encodings for processing extended input sequences. Studies such as Code LLaMA~\cite{codellama}, LLaMA2 Long~\cite{xiong2023effective}, and ScaledRoPE~\cite{liu2023scaling} implement a considerably larger base parameter for RoPE and subsequently fine-tune the model for the desired context length. In contrast, our approach supports a wide range of target lengths and is capable of exceeding a 2M sequence length. More recently, as extending context lengths beyond 128k through fine-tuning introduces significant computational demands, methods like LongLoRA~\cite{longlora} and PoSE~\cite{zhu2023pose} have been introduced to reduce these resource requirements. Our technique operates independently and is complementary to these resource-efficient fine-tuning methods.

\section{Conclusion}

This paper introduces LongRoPE, a technique that significantly increases the context window of LLMs to a remarkable 2048k, all while preserving their proficiency within the original, more limited window. The approach leverages two distinct types of non-uniformity in RoPE positional embeddings, discovered via efficient evolutionary search. This dual strategy achieves both an optimal starting point for subsequent fine-tuning and facilitates an 8$\times$ expansion of the context window in the absence of fine-tuning. Additionally, we present a staged extension methodology, whereby LLMs fine-tuned on 256k-length contexts are incrementally adapted to support 2048k context length without requiring further fine-tuning steps. Comprehensive evaluations demonstrate the robustness and efficacy of LongRoPE. We anticipate that the LongRoPE-2048k models will unlock novel applications requiring extended contexts and stimulate continued innovation in the field.

\section*{Broader Impacts} The research outlined in this paper aims to contribute to progress within Machine Learning. While our findings may have a range of implications for society, we have not identified any particular effects that warrant special emphasis in this context.

	\balance

\end{document}